\chapter{When timing noise and finite exposure preserve an inference}
\label{ch:virtual-timing}

A timing error need not be a small perturbation, and a finite exposure need not
behave like an instantaneous detector. Some large timing fluctuations preserve
the inference exactly. Other schedules produce apparent equilibrium violations
even when their mean lags are perfectly calibrated. This chapter identifies the
relevant contracts and gives exact controls for their failure.

\section{A universal clock criterion}
For positive integer settings $k=1,\ldots,K$, let $T_k\geq0$ be a random
observation time independent of the physical process. Define
\[
 \mathcal H_K=\left\{\left(\int_0^1x^k\,\mu(dx)\right)_{k=1}^K:
 \mu\text{ is a probability measure}\right\}.
\]
The set is closed and convex. A reversible stationary process, observed through
a fixed binary endpoint detector, has a correlation which is a positive mixture
of decaying exponentials. This standard spectral fact underlies the following
apparatus criterion~\cite{vlAldousFill2002,vlSchilling2012}.

\begin{theorem}[Clock compatibility]
The clock family preserves $\mathcal H_K$ for every finite stationary reversible
Markov chain and every fixed, noninvasive binary endpoint detector with independent
readout noise if and only if
\begin{equation}
 \left(\mathbb E e^{-\lambda T_k}\right)_{k=1}^K\in\mathcal H_K
 \quad\text{for every }\lambda\geq0.
 \label{eq:clock-criterion}
\end{equation}
\end{theorem}
\begin{proof}
Write the reversible correlation as
$C(t)=\sum_\ell b_\ell e^{-\lambda_\ell t}$, where $b_\ell\geq0$ and
$\sum_\ell b_\ell\leq1$. If~\eqref{eq:clock-criterion} holds, combine the
representing measures with weights $b_\ell$ and place the unused mass at zero.
All indices are positive, so the resulting measure represents the sampled
correlations. Conversely, a symmetric two-state chain with rates $\lambda/2$
and detector values $\pm1$ has exactly $C(t)=e^{-\lambda t}$. It tests each
scalar condition in~\eqref{eq:clock-criterion}.
\end{proof}

The theorem concerns a known clock family, not a finite numerical test over a
few relaxation rates. For an endpoint detector with calibrated stationary mean
$m$, its spectral atom $m^2$ at one is retained and the earlier mean-conditioned
cone applies. Finite signed differences are useful necessary checks; they are
not a complete finite moment-membership algorithm.

\section{Two safe families with different dimension implications}
If $T_k=X_1+\cdots+X_k$ for independent, identically distributed nonnegative
increments, then
\[
 \mathbb E e^{-\lambda T_k}=\left(\mathbb E e^{-\lambda X}\right)^k,
 \qquad
 \mathbb E e^{QT_k}=P^k,
 \qquad P=\mathbb E e^{QX}.
\]
The stochastic matrix $P$ acts on the same hidden state space. Thus arbitrary
common increment noise preserves both the reversible moment cone and the
classical dimension factorization. Neither exponential increments nor a hidden
rate cap is required.

A common random scale $T_k=kV$ also satisfies the clock criterion, by the
pushforward $x=e^{-\lambda V}$. It need not produce a common channel on the
same state space. For example, a two-state chain with decay $2^{-t}$ and a fair
scale choice $V\in\{1,2\}$ gives prepared occupation probabilities
\[
 p_k=\frac12+\frac{2^{-k}+4^{-k}}4.
\]
The $3\times3$ Hankel matrix $(p_{i+j})_{i,j=0}^2$ has determinant
$9/32768>0$. A two-state common-channel representation is impossible. A
three-state representation exists: two transient states survive each step with
probabilities $1/2$ and $1/4$, then enter an absorbing state; use equal initial
weights on the transient states and detector values $(1,1,1/2)$.
The pooled record includes a persistent scale variable. Preserving a
thermodynamic null therefore does not automatically preserve a dimension null.

\section{The actual duration belongs in the entropy bound}
Keep the full joint record consisting of the setting, the realized clock and
the physical trajectory. Its clock law is common to the actual generator $Q$
and its additive reversibilization $S$. The stationary path comparison gives
\[
 D(\mathbb P_Q\Vert\mathbb P_S)
 =\mathbb E[T]\,d(Q,S)\leq\frac{\mathbb E[T]}4\sigma,
 \qquad
 \mathbb E[T]=\sum_k w_k\mathbb E[T_k].
\]
Discarding the realized time only decreases relative entropy. Consequently,
for a certified record divergence $J$ from the compatible reversible cone,
\begin{equation}
 \sigma\geq\frac{4J}{\sum_kw_k\mathbb E[T_k]}.
 \label{eq:random-clock-entropy}
\end{equation}
The mean duration must be finite and positive; a calibrated upper bound on it
can be used conservatively. For common increments it equals
$\mathbb E[X]\sum_kw_kk$, and for common scale it equals
$\mathbb E[V]\sum_kw_kk$. This remains a stationary even-state classical
inference with the physical assumptions of the earlier entropy chapter.

Mean timing alone is insufficient. In an equilibrium two-state process with
$C(t)=2^{-t}$, take $T_1=1$ and let $T_2$ be zero with probability $4/5$ and ten
otherwise. The mean lags are one and two, but the correlations increase from
$1/2$ to $4097/5120$. Even stochastic ordering is insufficient for the full
moment cone: the coupled schedules $(1,1,2)$ and $(1,11,12)$, chosen with
probabilities $9/10$ and $1/10$, have means $(1,2,3)$ and pointwise ordered
lags, yet their second correlation difference is $-7171/40960$.
Both examples have zero physical entropy production.

\section{Finite-duration symmetric detectors}
Let a block last $L$ and have a bounded conditional binary mean
$F(X_{[0,L]})\in[-1,1]$. Assume $F$ is invariant under reversing the path within
the block, the detector is noninvasive, and different block outputs have
independent noise conditional on the physical paths. Use identical blocks on
$[0,L]$ and $[L+T,2L+T]$, where $T\geq0$ is the independent gap.

\begin{theorem}[Disjoint symmetric exposure]
For the reversible reference process, put
$g(x)=\mathbb E_S[F(X_{[0,L]})\mid X_0=x]$. Then
\begin{equation}
 \mathbb E_S[Y_{\rm left}Y_{\rm right}\mid T]
 =\langle g,e^{ST}g\rangle_\pi.
 \label{eq:block-positive-spectrum}
\end{equation}
Thus a compatible gap clock preserves the reversible moment cone. The
corresponding stationary entropy bound is
\begin{equation}
 \sigma\geq\frac{4J}{2L+\mathbb E[T]}.
 \label{eq:block-entropy}
\end{equation}
\end{theorem}
\begin{proof}
Stationary path reversal and symmetry of $F$ imply that its conditional mean
given the final block state is also $g$. Conditioning on the state at the end
of the first block and the beginning of the second, the Markov property yields
\eqref{eq:block-positive-spectrum}. Since $|g|\leq1$, its spectral measure has
mass at most one. The observed record uses a path of length $2L+T$; applying
the same path divergence comparison gives~\eqref{eq:block-entropy}.
\end{proof}

For a nonlinear block functional, the mean under $S$ can differ from the mean
under $Q$. The mean-unrestricted outer cone is safe: its agreement baseline is $1/2$,
and any conserved spectral mass is included. This does not assume that the
physical mean is zero. Measured mean information cannot simply be transferred
between the models. A symmetric
linear exposure,
\[
 F=\int_0^Lw(u)f(X_u)\,du,
 \qquad w\geq0,\quad\int_0^Lw(u)\,du=1,\quad w(u)=w(L-u),
\]
has mean $\pi f$ under either generator, so the calibrated-mean refinement
is valid. Rectangular exposure is included at the level of this bounded
conditional mean. Invasiveness, additional detector memory and conditional
noise dependence still require separate treatment.

\section{Three exposure controls that prevent false discoveries}
All three controls use an equilibrium two-state process with $C(t)=2^{-t}$.
A symmetric block channel reports a fair random choice of its two endpoint
signs. For length $L$ and start separation $D$, its correlation is exactly
\[
 C_{L}(D)=\frac{2\,2^{-D}+2^{-(D+L)}+2^{-|D-L|}}4.
\]
For disjoint blocks, $D=L+T$, this is
$((1+2^{-L})/2)^2 2^{-T}$, confirming the gap formula.

First, overlap can create a false revival. For $L=4$, start separations two
and four give $49/256$ and $289/1024$, an increase of $93/1024$.

Second, even disjoint blocks require the correct time coordinate. For $L=2$,
start separations two and four give $C_1=25/64$ and $C_2=25/256$. Assigning them
to the first and second normalized moments yields
$C_2-C_1^2=-225/4096<0$. The actual gaps are zero and two. A gap-zero correlation
has unknown spectral mass at most one and must not be assigned the normalization
of an ideal noiseless instantaneous detector.

Third, reversing the detector can reverse the sign test. The antisymmetric
channel $F=(X_{\rm end}-X_{\rm start})/2$, implemented through an independent
binary encoding, gives correlation $-9/128$ at length two and gap one.
Its agreement probability is $119/256<1/2$ despite equilibrium. The required
time-reversal symmetry of the block channel is absent.

These examples are exact synthetic population calculations, not laboratory
measurements. They show why the apparatus contract is part of the theorem.
The clock criterion is a direct consequence of established positive spectral
calculus; the present contribution is its integration with the inference
contracts, complete duration accounting and explicit failure controls.
